\section{BroadcastGuard system design}
Figure~\ref{fig:architecture} shows the three computational paths and the separate evidence pipeline. The educational attack modules deliberately use textbook RSA; secure file sharing uses standard authenticated encryption and never calls the textbook operation. This separation prevents a vulnerable research primitive from becoming a file-encryption dependency.

\begin{figure}[!ht]
\centering
\resizebox{\linewidth}{!}{\begin{tikzpicture}[
  x=1cm,y=1cm,
  box/.style={draw=usthblue,rounded corners=2pt,align=center,
              minimum width=2.75cm,minimum height=0.88cm,inner sep=3pt,font=\footnotesize},
  evidence/.style={draw=black!65,align=center,minimum width=8.6cm,
                   minimum height=0.70cm,inner sep=3pt,font=\footnotesize},
  flow/.style={-{Latex[length=2.2mm]},line width=0.65pt}
]
\node[box] (a1) at (0,2.6) {Shared bytes\\textbook builder, $e=3$};
\node[box] (a2) at (3.65,2.6) {Public observations\\$(n_i,e,c_i)$};
\node[box] (a3) at (7.3,2.6) {Eve: CRT and\\exact cube root};
\node[box] (a4) at (10.95,2.6) {Evaluator only:\\compare with $M$};
\draw[flow] (a1) -- (a2); \draw[flow] (a2) -- (a3); \draw[flow] (a3) -- (a4);

\node[box] (b1) at (0,1.35) {Same keys and bytes\\textbook / OAEP};
\node[box] (b2) at (3.65,1.35) {Two sets of public\\ciphertexts};
\node[box] (b3) at (7.3,1.35) {Same Eve method\\on each scheme};
\node[box] (b4) at (10.95,1.35) {Evaluator and\\valid OAEP recipients};
\draw[flow] (b1) -- (b2); \draw[flow] (b2) -- (b3); \draw[flow] (b3) -- (b4);

\node[box] (c1) at (0,0.1) {File bytes and verified\\recipient public keys};
\node[box] (c2) at (3.65,0.1) {AES-GCM once\\OAEP key wrap per ID};
\node[box] (c3) at (7.3,0.1) {Versioned JSON\\header as GCM AAD};
\node[box] (c4) at (10.95,0.1) {Matching recipient\\unwraps and decrypts};
\draw[flow] (c1) -- (c2); \draw[flow] (c2) -- (c3); \draw[flow] (c3) -- (c4);

\node[evidence] (ev) at (5.475,-1.15) {E1--E4 raw CSVs $\rightarrow$ read-back validation $\rightarrow$ descriptive tables and figures};
\draw[flow] (a4.south) |- (ev.east);
\draw[flow] (c3.south) |- (ev.north);
\end{tikzpicture}
}
\caption{Implemented data flows. Eve receives only public observations; original plaintext reaches the evaluator, not the attack engine. Defense compares two schemes under the same RSA keys. The hybrid path wraps one fresh AES key per recipient after encrypting the file once. The evidence pipeline validates archived experiment records before analysis.}
\label{fig:architecture}
\end{figure}

The \emph{scenario builder} converts input bytes to $M$, generates three RSA-2048/$e=3$ keypairs, checks $M<n_i$ and pairwise coprimality, and computes $c_i=\operatorname{pow}(M,3,n_i)$. Only public numbers survive in its observations. The \emph{attacker engine} accepts selected public observations, rejects duplicate or incompatible inputs, performs the explicit CRT computation in Eq.~\eqref{eq:crt}, and computes an integer cube root with an exactness flag. It is not given plaintext or a private key. The \emph{experiment evaluator} separately knows $M$ and reports whether a candidate matches it. The \emph{conditions analyzer} invokes these components for all seven nonempty subsets, including non-prefix pairs such as Bob--Charlie. It reports both the exact-root status and the evaluator-confirmed integer match, with the minimum successful interception count or an explicit no-success result.

The \emph{defense lab} uses the same three RSA-2048/$e=3$ keypairs and plaintext for textbook encryption and RSA-OAEP. It applies the same public CRT-and-root procedure to both sets of ciphertext integers; intended recipients also decrypt their OAEP ciphertexts with in-memory private keys. Holding the keys fixed isolates the scheme change within each paired trial. OAEP ciphertext integers are tested as attack inputs, but they are never mislabeled as textbook ciphertexts. Repeating OAEP encryption under one key gives a separate observation of ciphertext diversity.

The \emph{secure sharing} module uses RSA-2048/$e=65537$, not the educational $e=3$ configuration. Each file encryption creates a fresh 256-bit AES key and 96-bit GCM nonce. AES-GCM encrypts the file once; RSA-OAEP with SHA-256/MGF1 wraps the AES key independently for each of one to five authenticated recipient public keys. A canonical UTF-8 JSON header records format and algorithm identifiers, filename, original size, nonce, recipient IDs, public-key fingerprints and wrapped keys. Its serialized bytes serve as GCM associated authenticated data (AAD). The package contains Base64-encoded ciphertext with the 16-byte tag and no private or plaintext key material. The receiver validates structure and key fingerprint, unwraps its entry, reconstructs AAD and returns plaintext only after authenticated decryption and size verification~\cite{dworkin2007gcm}. Fingerprints label keys; they do not independently prove a person's identity.

The \emph{evaluation pipeline} stores individual E1--E4 records, a configuration manifest and explicit validation status. A read-back validator checks case coverage, paired-key consistency, timing identities and outcome invariants before statistical analysis. Figure generation reads the committed official CSVs; it does not rerun randomized cryptographic experiments. Source code, the run and reproduction instructions are available in the \href{https://github.com/Dung092005/Crypto-project}{repository}.
